% Propietario conservado: /Users/ruben/Documents/New project/output/EXTREMA_MEDIA_RAZON_RECIPROCIDAD_ES_EN_20260918/ARTICULO/ES/source/libro/05_levantamiento.tex
\chapter{Relèvement unitaire du raffinement et transport des lecteurs}
\label{x_reciprocidad:lib05:levantamiento}

La bijection par les restes de la \cref{x_reciprocidad:lib03:teo-residuo} et le
bilan bilatéral de la \cref{x_reciprocidad:lib04:teo-balance} admettent une
composition explicite. Un triplet arithmétique admissible, qui contient
un couple et sa retenue, détermine un couple raffiné. La
linéarisation de cette bijection sur des bases orthonormées fournit
un transport unitaire ; deux transports de ce type peuvent agir
comme les deux réalisations du bilan bilatéral.

La réalisation est établie après les objets engendrés par APP--TRIT--TPK. Son domaine combinatoire comprend les paires et les retenues qui satisfont les conditions écrites ci-dessous. La norme conservée par sa linéarisation est la somme des carrés des amplitudes sur ces étiquettes ; la conservation arithmétique de \(L_c(x,y)=(Bx-c,By+c)\) reste la somme des deux coordonnées après remise à l'échelle. La composition relie les deux conservations au moyen d'une application concrète.

\section{Bijection des états arithmétiques admissibles}
\label{x_reciprocidad:lib05:biyeccion}

Soient \(B\geq2\) et \(M\geq1\) des entiers. Nous définissons
\begin{align}
 \mathcal D_{B,M}
 &=\bigl\{(x,y,c)\in\mathbb Z_{\geq0}^2
                 \times\{0,\ldots,B-1\}:x+y=M,\ c\leq Bx\bigr\},
 \label{x_reciprocidad:lib05:dominio}\\
 \mathcal E_{BM}
 &=\bigl\{(X,Y)\in\mathbb Z_{\geq0}^2:X+Y=BM\bigr\}.
 \label{x_reciprocidad:lib05:codominio}
\end{align}
La condition \(c\leq Bx\) maintient la non-négativité de la
première coordonnée de la sortie. Pour \(x=0\), elle n’admet que
\(c=0\) ; pour \(x\geq1\), elle admet tout l’alphabet canonique.

\begin{theorem}[Correspondance exacte des étiquettes]
\label{x_reciprocidad:lib05:teo-biyeccion}
L'application
\begin{equation}
 \ell_{B,M}:\mathcal D_{B,M}\longrightarrow\mathcal E_{BM},
 \qquad (x,y,c)\longmapsto(Bx-c,By+c)
 \label{x_reciprocidad:lib05:mapa-etiquetas}
\end{equation}
est bijective. Les deux ensembles ont pour cardinal \(BM+1\). Son inverse
est donné par
\begin{equation}
 c=Y\bmod B,\qquad y=\frac{Y-c}{B},\qquad x=\frac{X+c}{B}.
 \label{x_reciprocidad:lib05:inversa-etiquetas}
\end{equation}
\end{theorem}

\begin{proof}
La condition de domaine et l'identité \(Bx-c+By+c=BM\) assurent que l'image appartient à \(\mathcal E_{BM}\). Étant donnée une paire de cet ensemble, le résidu \(c\) est l'unique élément de l'alphabet qui satisfait \(Y\equiv c\pmod B\). Comme \(X+Y\) est divisible par \(B\), \(X+c\) l'est également. La division euclidienne de \(Y\) donne \(y\geq0\), et \(X+c\geq0\) donne \(x\geq0\). Si \(c>0\), \(x=0\) exigerait \(X=-c<0\), de sorte que l'admissibilité est conservée. La somme des antécédents est \(M\), et la substitution vérifie l'image originale. L'unicité du résidu et des quotients démontre la bijection.

Il y a \(M+1\) étiquettes avec \(c=0\) et \(M\) pour chaque retenue
positive ; le cardinal est donc
\((M+1)+(B-1)M=BM+1\), égal à celui de la sortie.
\end{proof}

La preuve explicite l’extrémité de frontière du domaine, condition
nécessaire à l’égalité des cardinaux. La nouvelle étiquette de retenue
est une partie typée de l’entrée du raffinement.

\section{Unitarité sur les amplitudes des étiquettes}
\label{x_reciprocidad:lib05:unitariedad}

Prenons les espaces de Hilbert réels ou complexes munis de la mesure de comptage
\begin{equation}
 \mathscr H_D=\ell^2(\mathcal D_{B,M}),\qquad
 \mathscr H_E=\ell^2(\mathcal E_{BM}),
 \label{x_reciprocidad:lib05:hilbert}
\end{equation}
et leurs bases orthonormées étiquetées. Nous définissons linéairement
\begin{equation}
 J_{B,M}|x,y,c\rangle=|Bx-c,By+c\rangle.
 \label{x_reciprocidad:lib05:j}
\end{equation}
On écrira \(J\) lorsque les paramètres seront fixés.

\begin{proposition}[Relèvement unitaire]
\label{x_reciprocidad:lib05:prop-unitaria}
L’opérateur \(J:\mathscr H_D\to\mathscr H_E\) est unitaire :
\(J^*J=I_D\), \(JJ^*=I_E\). L’adjoint retrouve chaque étiquette
au moyen de \eqref{x_reciprocidad:lib05:inversa-etiquetas}.
\end{proposition}

\begin{proof}
La \cref{x_reciprocidad:lib05:teo-biyeccion} assure que \(J\) transforme
bijectivement deux bases orthonormées complètes. Pour
\(\psi=\sum_{d\in\mathcal D}\psi_d|d\rangle\),
\[
 \|J\psi\|^2
 =\sum_{e\in\mathcal E}|\psi_{\ell_{B,M}^{-1}(e)}|^2
 =\sum_{d\in\mathcal D}|\psi_d|^2=\|\psi\|^2.
\]
La surjectivité prouve les deux identités unitaires et la
bijection inverse décrit l’adjoint sur chaque vecteur de base.
\end{proof}

L’étiquette \(|73,27,1\rangle\) a norme un ; les entiers
\(73,27,1\) décrivent ses valeurs arithmétiques, non ses amplitudes.
La construction représente un traitement réversible d’étiquettes
ou d’amplitudes. Une réalisation physique concrète conserve en outre
le protocole qui prépare et transporte ces amplitudes.

\section{Lecteurs diagonaux de proportion et de retenue}
\label{x_reciprocidad:lib05:diagonales}

Pour une fonction réelle \(h\) sur un ensemble fini, notons
\(M_h\) son opérateur diagonal : chaque vecteur de base est multiplié
par la valeur de \(h\) en son étiquette. La bijection vérifie
\begin{equation}
 J^*M_hJ=M_{h\circ\ell_{B,M}}.
 \label{x_reciprocidad:lib05:lector-diagonal}
\end{equation}
Cela se démontre en agissant sur tous les vecteurs de la base.

Définissons en sortie
\begin{equation}
 f_E(X,Y)=\frac{X}{BM},\quad
 g_E(X,Y)=\frac{Y}{BM},\quad
 \gamma_E(X,Y)=Y\bmod B,
 \label{x_reciprocidad:lib05:lectores-salida}
\end{equation}
et en entrée \(f_D=x/M\), \(g_D=y/M\), \(\gamma_D=c\). L'identité \eqref{x_reciprocidad:lib05:lector-diagonal} fournit
\begin{align}
 J^*M_{f_E}J&=M_{\,x/M-c/(BM)},\nonumber\\
 J^*M_{g_E}J&=M_{\,y/M+c/(BM)},\nonumber\\
 J^*M_{\gamma_E}J&=M_c.
 \label{x_reciprocidad:lib05:transporte-fracciones}
\end{align}
Les deux premiers lecteurs transportent le transfert d’unité.
Comme \(f_E+g_E=1\), leur somme est \(I_E\) et leur compression est
\(I_D\). Le troisième lecteur conserve exactement la retenue comme
observable entier, avec sa représentation résiduelle dans la sortie.

\section{Chemins complets à profondeur arbitraire}
\label{x_reciprocidad:lib05:rutas}

Soit \(\mathcal P_n(B,M)\) l'ensemble des chemins de longueur \(n\) qui partent de toutes les paires non négatives de somme \(M\) et, à chaque étape, satisfont les conditions de \eqref{x_reciprocidad:lib05:dominio} pour la somme correspondante. Un chemin conserve sa paire initiale et ses \(n\) retenues admissibles. Soit \(\Phi_n\) l'application qui attribue à chaque chemin sa paire finale.

\begin{theorem}[Récupération de routes à toute profondeur finie]
\label{x_reciprocidad:lib05:teo-rutas}
Pour tout \(n\geq0\),
\begin{equation}
 \Phi_n:\mathcal P_n(B,M)\longrightarrow\mathcal E_{B^nM}
 \quad\hbox{est bijective},\qquad
 |\mathcal P_n(B,M)|=B^nM+1.
 \label{x_reciprocidad:lib05:biyeccion-rutas}
\end{equation}
Sa linéarisation
\begin{equation}
 J_n:\ell^2(\mathcal P_n(B,M))\longrightarrow
          \ell^2(\mathcal E_{B^nM}),\qquad
 J_n|\pi\rangle=|\Phi_n(\pi)\rangle
 \label{x_reciprocidad:lib05:jn}
\end{equation}
est unitaire.
\end{theorem}

\begin{proof}
À partir d’une paire de somme \(B^nM\), le résidu de
\eqref{x_reciprocidad:lib05:inversa-etiquetas} retrouve la dernière retenue et une
unique paire précédente de somme \(B^{n-1}M\). Répéter \(n\) fois
produit une unique route dont l’entrée a pour somme \(M\). La composition
directe retrouve l’extrémité originale, ce qui prouve la bijection et
l’inverse. Le dénombrement s’obtient également à partir des entrées : les
\(M\) paires avec \(x_0\geq1\) admettent chacune \(B^n\) mots,
tandis que \((0,M)\) n’admet que le mot entièrement
nul. Le total est \(MB^n+1\). La transformation bijective des
bases orthonormales démontre l’unitarité de \(J_n\).
\end{proof}

La profondeur étant connue, les lecteurs sur le couple final sont
\begin{equation}
 c_j=\left\lfloor\frac{Y}{B^{n-j}}\right\rfloor\bmod B,
 \quad y_0=\left\lfloor\frac{Y}{B^n}\right\rfloor,
 \quad x_0=M-y_0,
 \qquad 1\leq j\leq n.
 \label{x_reciprocidad:lib05:lectores-ruta}
\end{equation}
Ces fonctions sont diagonales et récupèrent la mémoire complète des
opérations finies. La compatibilité avec le prolongement est
\begin{equation}
 \Phi_{n+1}(\pi,c)=L_c(\Phi_n(\pi)).
 \label{x_reciprocidad:lib05:compatibilidad-prolongacion}
\end{equation}
Lors du passage de \(n\) à \(n+1\), on incorpore la nouvelle étiquette
admissible de retenue. Les espaces augmentent leur dimension et
l’étiquette de frontière n’admet que \(c=0\) ; le domaine de
l’étape suivante conserve donc l’admissibilité, au lieu de la remplacer
par un produit sans restriction avec \(B\) chiffres.

\section{Entrée fixée et image exacte}
\label{x_reciprocidad:lib05:entrada-fija}

Si l'on fixe une paire initiale \((x_0,y_0)\), avec \(x_0\geq1\), il existe \(B^n\) chemins et leur image est exactement
\begin{equation}
 \bigl\{(B^nx_0-C,B^ny_0+C):0\leq C\leq B^n-1\bigr\}.
 \label{x_reciprocidad:lib05:imagen-fija}
\end{equation}
La \cref{x_reciprocidad:lib03:cor-palabra} démontre que les \(B^n\) cumuls sont distincts et récupèrent les mots respectifs. La linéarisation est unitaire vers le sous-espace engendré par cette image. Si \(x_0=0\), l'image est constituée d'une seule étiquette.

Les cylindres de la \cref{x_reciprocidad:lib03:teo-cilindros} correspondent à une
entrée fixée. Le domaine complet de
\eqref{x_reciprocidad:lib05:biyeccion-rutas} réunit toutes les paires de somme \(M\).
Les deux constructions possèdent des inverses explicites, mais leurs images
et leurs cardinaux sont différents. La profondeur arbitraire exprime la
validité de la preuve pour tout \(n\) ; la limite de la lecture et
ses frontières restent celles établies dans la
\cref{x_reciprocidad:lib03:limite}.

\section{Couplage bilatéral du raffinement}
\label{x_reciprocidad:lib05:bilateral}

Soit \(R\) une permutation préalablement déclarée de
\(\mathcal D_{B,M}\), linéarisée unitairement dans
\(\mathscr H_D\). Nous définissons les deux transports
\begin{equation}
 \mathcal B=J,\qquad \mathcal C=JR:
            \mathscr H_D\longrightarrow\mathscr H_E.
 \label{x_reciprocidad:lib05:dos-transportes}
\end{equation}
La permutation fait partie de la spécification de la procédure et est fixée avant l'évaluation d'une lecture de sortie. Pour \(a>0\), \(p=a^2+a+1\), \(N=(2a+1)p\),
\begin{equation}
 Q_-=J(aI-R),\qquad Q_+=J((a+1)I+R).
 \label{x_reciprocidad:lib05:q-refinamiento}
\end{equation}

\begin{proposition}[Distribution et récupération des amplitudes raffinées]
\label{x_reciprocidad:lib05:prop-bilateral}
Les opérateurs \eqref{x_reciprocidad:lib05:q-refinamiento} satisfont
\begin{equation}
 (a+1)Q_-^*Q_-+aQ_+^*Q_+=NI_D,
 \label{x_reciprocidad:lib05:balance}
\end{equation}
et l’application
\begin{equation}
 W_a\psi=\frac1{\sqrt N}
    \bigl(\sqrt{a+1}\,Q_-\psi,\sqrt a\,Q_+\psi\bigr)
 \label{x_reciprocidad:lib05:w}
\end{equation}
est isométrique. Pour les canaux non normalisés \(u=Q_-\psi\), \(v=Q_+\psi\),
\begin{equation}
 J\psi=\frac{u+v}{2a+1},\qquad
 \psi=J^*\frac{u+v}{2a+1}.
 \label{x_reciprocidad:lib05:inversa-bilateral}
\end{equation}
\end{proposition}

\begin{proof}
Les deux applications \(\mathcal B,\mathcal C\) sont unitaires.
Les termes croisés des gramiens de \(Q_-\) et de \(Q_+\) ont pour
coefficients \(-a\) et \(a+1\), respectivement. Les pondérations
du bilan les annulent, tandis que les termes diagonaux ont pour somme
\(N\), comme dans la \cref{x_reciprocidad:lib04:teo-balance}. La normalisation
démontre l’isométrie. La somme des canaux est
\((2a+1)J\psi\), et l’unitarité de \(J\) donne l’inverse.
\end{proof}

À partir des canaux normalisés, on récupère directement l'entrée par \(W_a^*\), avec une norme d'inversion égale à un ; cet adjoint incorpore déjà \(J^*\). La formule alternative \eqref{x_reciprocidad:lib05:inversa-bilateral} sépare la récupération de \(J\psi\) par somme de l'application de \(J^*\), dont l'action sur les étiquettes se calcule au moyen des résidus. Le raffinement arithmétique réversible et le transport bilatéral des amplitudes sont ainsi composés.

\begin{proposition}[Compatibilité dans l’espace raffiné]
\label{x_reciprocidad:lib05:prop-compatibilidad}
Une paire non normalisée \((u,v)\) appartient à l’image de
\(\psi\mapsto(Q_-\psi,Q_+\psi)\) si et seulement si
\begin{equation}
 \bigl[(a+1)I_E+JRJ^*\bigr]u
 =\bigl[aI_E-JRJ^*\bigr]v.
 \label{x_reciprocidad:lib05:compatibilidad}
\end{equation}
\end{proposition}

\begin{proof}
L’égalité équivaut à
\(JRJ^*(u+v)=av-(a+1)u\). Si elle est satisfaite, définissons
\(\psi=J^*(u+v)/(2a+1)\). Lors de la substitution dans
\eqref{x_reciprocidad:lib05:q-refinamiento}, cette identité donne
\(Q_-\psi=u\) et \(Q_+\psi=v\). La nécessité résulte de
l’insertion des canaux véritables. Toutes les compositions agissent
dans \(\mathscr H_E\), de sorte que leurs types sont définis.
\end{proof}

\section{Loi des observables sur les deux canaux}
\label{x_reciprocidad:lib05:ley-lectores}

\begin{theorem}[Compression bilatérale d'un lecteur raffiné]
\label{x_reciprocidad:lib05:teo-lector}
Pour \(F=F^*\) sur \(\mathscr H_E\), écrivons
\(F_J=J^*FJ\). Alors
\begin{equation}
 W_a^*\operatorname{diag}(F,F)W_a
 =\frac{a(a+1)F_J+R^*F_JR}{p}.
 \label{x_reciprocidad:lib05:lector-bilateral}
\end{equation}
Le lecteur \(F_J\) est conservé exactement si et seulement s'il commute avec \(R\).
\end{theorem}

\begin{proof}
Développer \((a+1)Q_-^*FQ_-+aQ_+^*FQ_+\) annule les
termes mixtes. Le coefficient de \(\mathcal B^*F\mathcal B\)
est \(a(a+1)(2a+1)\) ; celui de
\(\mathcal C^*F\mathcal C\) est \(2a+1\). Diviser par
\(N\), substituer \(\mathcal C=JR\) et utiliser
\(F_J=J^*FJ\) prouve la formule. L’égalité avec \(F_J\)
équivaut à \(R^*F_JR=F_J\), ou à \([F_J,R]=0\) par
unitarité.
\end{proof}

Dans le cas \(a=8\), les poids sont \(72/73\) et \(1/73\). Le lecteur identité récupère la norme. Un lecteur non constant peut changer même lorsque la norme est conservée ; sa variation est calculée par \eqref{x_reciprocidad:lib05:lector-bilateral}. L'opérateur qui conserve exactement un lecteur original \(A\) dans l'image complète est \(W_aAW_a^*\), dont l'identité d'image est \(W_aW_a^*\). Lire \(\operatorname{diag}(F,F)\) correspond à la compression spécifique du théorème et ne remplace pas automatiquement cet observable transporté.

\section{Exemple arithmétique de dimension mille un}
\label{x_reciprocidad:lib05:ejemplo}

Prenons \(B=10\), \(M=100\). Les espaces d’étiquettes sont de
dimension \(1001\), et
\begin{equation}
 d_1=(73,27,1)\longmapsto e_1=(729,271).
 \label{x_reciprocidad:lib05:ejemplo-etiqueta}
\end{equation}
Les lecteurs diagonaux donnent
\begin{equation}
 f_E(e_1)=\frac{729}{1000}
 =\frac{73}{100}-\frac1{1000},\qquad
 g_E(e_1)=\frac{271}{1000},\qquad \gamma_E(e_1)=1.
 \label{x_reciprocidad:lib05:ejemplo-lectores}
\end{equation}
Déclarons \(R\) comme la transposition des étiquettes
admissibles \(d_1=(73,27,1)\) et \(d_2=(73,27,2)\), en laissant
les autres fixes. La seconde image est \(e_2=(728,272)\).
La règle de transposition a été fixée sur les étiquettes, avant
toute comparaison avec une constante.

Pour \(a=8\) et \(\psi=|d_1\rangle\),
\begin{equation}
 Q_-\psi=8|e_1\rangle-|e_2\rangle,\qquad
 Q_+\psi=9|e_1\rangle+|e_2\rangle.
 \label{x_reciprocidad:lib05:ejemplo-canales}
\end{equation}
Les carrés des normes sont \(65\) et \(82\), et
\begin{equation}
 9\cdot65+8\cdot82=585+656=1241=17\cdot73.
 \label{x_reciprocidad:lib05:ejemplo-balance}
\end{equation}
La somme des canaux est \(17|e_1\rangle\). L’inverse
\eqref{x_reciprocidad:lib05:inversa-bilateral} et le reste récupèrent
\(|d_1\rangle\), y compris la retenue un.

Si l’on évalue \(F=M_{f_E}\) diagonalement dans les deux canaux
normalisés, la \cref{x_reciprocidad:lib05:teo-lector} donne
\begin{equation}
 \begin{split}
 \langle W_8\psi,\operatorname{diag}(F,F)W_8\psi\rangle
 &=\frac{72\cdot729+728}{73\cdot1000}\\
 &=\frac{729}{1000}-\frac1{73000}.
 \end{split}
 \label{x_reciprocidad:lib05:ejemplo-fraccion}
\end{equation}
La lecture complémentaire est un moins cette quantité. Pour le
lecteur diagonal de retenue, on obtient
\begin{equation}
 \frac{72\cdot1+2}{73}=\frac{74}{73}.
 \label{x_reciprocidad:lib05:ejemplo-acarreo}
\end{equation}
La retenue originale reste récupérable exactement :
l’observable \(W_8M_cW_8^*\) a sur \(W_8|d_1\rangle\)
la valeur propre un. La valeur \(74/73\) est la réponse du
lecteur diagonal appliqué aux deux canaux, dont le transport relatif
est spécifié par \eqref{x_reciprocidad:lib05:lector-bilateral}. L’exemple
distingue conservation de l’information et variation d’une lecture
particulière.

\section{Articulation avec le registre engendré et domaines d'application}
\label{x_reciprocidad:lib05:dominios}

L’unitaire \(J\) agit sur \(BM+1\) étiquettes, tandis que
la réalisation de \(K\) a douze coordonnées. Dans l’exemple,
\(1001\) compte les paires raffinées de somme mille, et douze compte
les fenêtres du registre produit. Les deux dimensions conservent leur
origine combinatoire spécifique.

Une composition particulière entre ces deux espaces porteurs peut se formuler au moyen d'un sélecteur matériel de douze chemins distincts, avec leur provenance et leurs opérateurs. Si ce sélecteur fournit une isométrie \(J_{12}:\mathbb R^{12}\to\ell^2(\mathcal D_{B,M})\), alors \(JJ_{12}\) est isométrique, et l'entrelacement des lecteurs se vérifie sur celle-ci. Telle est l'hypothèse de cette composition supplémentaire, et non une identification implicite par le nombre de coordonnées. Les applications bilatérales déjà définies directement sur \(K\) dans la \cref{x_reciprocidad:lib04:bilateral} conservent leurs domaines complets sans utiliser ce sélecteur supplémentaire.

Le résultat de cette section réunit une bijection arithmétique, son
unitaire, des lecteurs diagonaux entrelacés, la récupération de routes à
toute profondeur finie et une composition bilatérale explicite.
Le domaine des routes est celui défini dans
\eqref{x_reciprocidad:lib05:biyeccion-rutas}. Les applications aux histoires
filtrées par TPK conservent leurs règles de survie et de
prolongement ; la correspondance des étiquettes prouve ici le
transport déclaré, sans remplacer ces opérateurs par un domaine
arithmétique sans restriction.
